\documentclass[conference]{IEEEtran}
\usepackage[utf8]{inputenc}
\usepackage{amsmath,amssymb,amsfonts}
\usepackage{algorithmic}
\usepackage{graphicx}
\usepackage{textcomp}
\usepackage{xcolor}
\usepackage{booktabs}
\usepackage{multirow}
\usepackage{url}
\usepackage{cite}

\begin{document}

\title{A High-Concurrency Peer-to-Peer Campus Task Dispatcher with ACID-Compliant Escrow Isolation and Pessimistic Locking}

\author{
    \IEEEauthorblockN{Rajveer Modi}
    \IEEEauthorblockA{\textit{Department of Computer Science and Engineering} \\
    \textit{Vellore Institute of Technology}\\
    Vellore, Tamil Nadu, India \\
    Email: rajveer.modi2025@vitstudent.ac.in}
}

\maketitle

\begin{abstract}
Decentralized peer-to-peer (P2P) micro-economies within university campuses offer substantial utility for micro-errands such as academic document dispatch, lab courier services, and food deliveries. However, naive database architectures in campus dispatch systems suffer from financial race conditions, including double-spending of bounties, phantom runner claims, and table-scan bottlenecks under high concurrent access. In this paper, we architect, implement, and empirically evaluate an ACID-compliant peer-to-peer campus task dispatching engine backed by MySQL 8.0 InnoDB and FastAPI. The underlying relational schema is formally proved to satisfy Boyce-Codd Normal Form (BCNF) through lossless-join decomposition and complete functional dependency preservation. To eliminate financial vulnerabilities without distributed consensus overhead, we design a centralized pessimistic row-level locking mechanism (\texttt{SELECT ... FOR UPDATE}) coupled with an isolated escrow transaction ledger and cryptographic 4-digit PIN verification. Empirical benchmark results on a 10,000-task workload demonstrate that our composite B-Tree indexing scheme reduces mean query latency by 33.1\% for open marketplace feeds (from 6.8837~ms to 4.6063~ms) and by 91.1\% ($11.2\times$ acceleration) for runner bid lookups (from 4.7457~ms to 0.4226~ms). Furthermore, under high-contention multi-threaded race conditions (including 10-thread parallel payout attacks), the engine maintains zero escrow leakage and achieves 100\% mathematical balance invariance with a mean contention resolution latency of 14.97~ms.
\end{abstract}

\begin{IEEEkeywords}
Peer-to-Peer Dispatch, ACID Transactions, Escrow Isolation, Pessimistic Locking, BCNF Normalization, Database Indexing, Concurrency Control.
\end{IEEEkeywords}

\section{Introduction}
Micro-errand fulfillment within dense university campus environments presents unique distributed systems and database challenges \cite{gray1992transaction}. Campus participants frequently exchange peer-to-peer delivery requests across geographically separated facilities (e.g., student hostels, academic towers, and research laboratories). While mobile platforms have democratized gig economies \cite{chen2019ridesharing}, conventional web frameworks often delegate balance deductions and task assignments to application-level logic without strict database transaction isolation. 

Under peak conditions (e.g., examination submission deadlines or rainy weather surges), concurrent claims on high-value errands expose critical race conditions:
\begin{enumerate}
    \item \textbf{Double-Spending of Escrow Funds:} Multiple threads concurrently invoking payout endpoints can execute duplicate wallet credits if the task completion and balance update are not isolated within an exclusive lock boundary.
    \item \textbf{Simultaneous Task Claiming:} Multiple competing student runners attempting to accept an errand simultaneously can result in conflicting task states or orphaned runner allocations.
    \item \textbf{Marketplace Manipulation \& Self-Bidding:} Requesters bidding on their own errands to artificially manipulate reputation metrics.
    \item \textbf{Query Latency Degradation:} High-volume task tables experiencing full table scans on status-driven feeds.
\end{enumerate}

To address these vulnerabilities, we design an enterprise-grade, relational peer-to-peer task dispatching platform. Our key contributions are:
\begin{itemize}
    \item \textbf{Formal Schema Normalization:} Complete algebraic derivation demonstrating that the database entities (\texttt{Students}, \texttt{Tasks}, \texttt{Bids}, and \texttt{Escrow\_Transactions}) satisfy strict Boyce-Codd Normal Form (BCNF) with guaranteed lossless join and dependency preservation.
    \item \textbf{Pessimistic Escrow Isolation:} Stored procedures utilizing InnoDB exclusive row-level locks (\texttt{FOR UPDATE}) and transactional state transitions enforcing absolute financial invariants.
    \item \textbf{Event-Driven Active Integrity:} MySQL triggers preventing self-bidding at the engine layer and automatically adjudicating competing runner bids upon task assignment.
    \item \textbf{Comprehensive Empirical Benchmarking:} Rigorous stress-testing of composite B-Tree index structures against full table scans on a 10,000-entity dataset, alongside multi-threaded race condition simulations.
\end{itemize}

\section{Relational Schema Architecture \& BCNF Normalization}
The relational database comprises four core entities engineered under the InnoDB storage engine to provide strict ACID guarantees.

\subsection{Entity Formulations}
\begin{itemize}
    \item $\text{Students}(\underline{\text{Reg\_No}}, \text{Name}, \text{Email}, \text{Trust\_Rating}, \text{Balance}, \text{Created\_At})$: Represents verified participants identified by institutional registration numbers.
    \item $\text{Tasks}(\underline{\text{Task\_ID}}, \text{Requester\_ID}, \text{Runner\_ID}, \text{Pickup\_Location}, \text{Drop\_Location}, \text{Bounty}, \text{Status}, \text{Verification\_Pin}, \text{Created\_At})$: Encapsulates delivery requests and routing coordinates.
    \item $\text{Bids}(\underline{\text{Bid\_ID}}, \text{Task\_ID}, \text{Runner\_ID}, \text{Offered\_Time}, \text{Bid\_Status}, \text{Created\_At})$: Captures runner bids for available tasks.
    \item $\text{Escrow\_Transactions}(\underline{\text{Txn\_ID}}, \text{Task\_ID}, \text{Amount}, \text{Escrow\_State}, \text{Released\_At}, \text{Updated\_At})$: Isolates held funds from student balances.
\end{itemize}

\subsection{Formal BCNF Proof}
A relation $R$ is in Boyce-Codd Normal Form (BCNF) if and only if for every non-trivial functional dependency $X \to Y \in F^+$, the determinant $X$ is a superkey of $R$ \cite{codd1974further}.

\subsubsection{Students Relation}
The functional dependency set is:
\begin{align*}
    F_1 = \{ & \text{Reg\_No} \to \{\text{Name}, \text{Email}, \text{Trust\_Rating}, \text{Balance}, \text{Created\_At}\}, \\
             & \text{Email} \to \{\text{Reg\_No}, \text{Name}, \text{Trust\_Rating}, \text{Balance}, \text{Created\_At}\} \}
\end{align*}
The attribute closures are $\{\text{Reg\_No}\}^+ = R_1$ and $\{\text{Email}\}^+ = R_1$. The candidate keys are $\mathcal{K}_1 = \{\{\text{Reg\_No}\}, \{\text{Email}\}\}$. Since both determinants are candidate keys, every determinant is a superkey. Therefore, $\text{Students}$ is in BCNF.

\subsubsection{Tasks Relation}
The functional dependency set is:
\begin{equation*}
    F_2 = \{ \text{Task\_ID} \to R_2 \setminus \{\text{Task\_ID}\} \}
\end{equation*}
The candidate key is $\mathcal{K}_2 = \{\{\text{Task\_ID}\}\}$. The determinant $\text{Task\_ID}$ is a superkey. Thus, $\text{Tasks}$ is in BCNF.

\subsubsection{Bids Relation}
The functional dependency set is:
\begin{align*}
    F_3 = \{ & \text{Bid\_ID} \to R_3 \setminus \{\text{Bid\_ID}\}, \\
             & \{\text{Task\_ID}, \text{Runner\_ID}\} \to R_3 \setminus \{\text{Task\_ID}, \text{Runner\_ID}\} \}
\end{align*}
The candidate keys are $\mathcal{K}_3 = \{\{\text{Bid\_ID}\}, \{\text{Task\_ID}, \text{Runner\_ID}\}\}$. Both determinants are candidate keys and superkeys. Therefore, $\text{Bids}$ is in BCNF.

\subsubsection{Escrow\_Transactions Relation}
The functional dependency set is:
\begin{align*}
    F_4 = \{ & \text{Txn\_ID} \to R_4 \setminus \{\text{Txn\_ID}\}, \\
             & \text{Task\_ID} \to R_4 \setminus \{\text{Task\_ID}\} \}
\end{align*}
Enforced by the candidate key constraint $\text{UNIQUE}(\text{Task\_ID})$, the candidate keys are $\mathcal{K}_4 = \{\{\text{Txn\_ID}\}, \{\text{Task\_ID}\}\}$. Both determinants are superkeys. Thus, $\text{Escrow\_Transactions}$ is in BCNF.

\subsection{Lossless-Join Decomposition and Dependency Preservation}
Let $U = R_1 \cup R_2 \cup R_3 \cup R_4$. By Heath's Theorem \cite{heath1971unacceptable}:
\begin{equation}
    (R_1 \cap R_2) = \{\text{Reg\_No}\} \to R_1 \implies R_1 \bowtie R_2 \text{ is lossless.}
\end{equation}
\begin{equation}
    (R_1 \cup R_2) \cap R_4 = \{\text{Task\_ID}\} \to R_4 \implies (R_1 \bowtie R_2) \bowtie R_4 \text{ is lossless.}
\end{equation}
\begin{equation}
    (R_1 \cup R_2 \cup R_4) \cap R_3 = \{\text{Task\_ID}, \text{Runner\_ID}\} \to R_3 \implies \text{Lossless.}
\end{equation}
Furthermore, because each functional dependency in $F$ is defined exclusively on attributes within its respective relation schema ($\bigcup \pi_{R_i}(F) \equiv F^+$), all dependencies are fully preserved without cross-relation join evaluation.

\section{Concurrency Control \& Pessimistic Locking Engine}
Financial transactions in peer-to-peer errand systems require strict serializability \cite{bernstein1987concurrency}. Optimistic concurrency control (OCC) incurs high abort rates under heavy write contention on single task records. Consequently, we implement a centralized pessimistic locking protocol using MySQL InnoDB row-level exclusive locks ($X$-locks).

\subsection{Escrow Lifecycle State Machine}
Fig.~\ref{fig:escrow_fsm} illustrates the state transitions governing task bounties. When a requester posts a task via \texttt{sp\_post\_task\_with\_escrow}, the user's balance row is locked with \texttt{SELECT ... FOR UPDATE}. Upon verifying $\text{Balance} \ge \text{Bounty}$, the bounty is deducted from the student's ledger and credited to an escrow record in the \texttt{HELD} state within a single ACID transaction.

\begin{figure}[htbp]
\centering
\fbox{
\begin{minipage}{0.9\columnwidth}
\centering
\vspace{4pt}
$\mathbf{Created} \xrightarrow[\text{Lock Requester}]{\text{sp\_post\_task\_with\_escrow}} \mathbf{HELD}$ \\
\vspace{6pt}
$\mathbf{HELD} \xrightarrow[\text{Valid PIN + Runner Match}]{\text{sp\_complete\_task\_handover}} \mathbf{RELEASED} \implies \text{Credit Runner}$ \\
\vspace{6pt}
$\mathbf{HELD} \xrightarrow[\text{Timeout / Cancellation}]{\text{sp\_refund\_abandoned\_task}} \mathbf{REFUNDED} \implies \text{Refund Requester}$
\vspace{4pt}
\end{minipage}
}
\caption{Escrow Transaction State Transition Machine.}
\label{fig:escrow_fsm}
\end{figure}

\subsection{Handover Verification Algorithm}
The payout handover is implemented in \texttt{sp\_complete\_task\_handover} as specified in Algorithm~\ref{alg:handover}.

\begin{algorithm}[htbp]
\caption{Pessimistic Task Handover and Escrow Release}
\label{alg:handover}
\begin{algorithmic}[1]
\REQUIRE $p\_task\_id$, $p\_runner\_id$, $p\_pin$
\STATE $\text{START TRANSACTION}$
\STATE $\text{SELECT Verification\_Pin, Status, Bounty, Runner\_ID}$
\STATE \quad $\text{INTO } v\_pin, v\_status, v\_bounty, v\_runner$
\STATE \quad $\text{FROM Tasks WHERE Task\_ID} = p\_task\_id \text{ \textbf{FOR UPDATE}};$
\IF{$v\_status \notin (\text{'assigned'}, \text{'in\_transit'})$}
    \STATE $\text{SIGNAL SQLSTATE '45000'} \implies \text{Task Not Active};$
\ENDIF
\IF{$v\_runner \ne p\_runner\_id$}
    \STATE $\text{SIGNAL SQLSTATE '45000'} \implies \text{Unauthorized Runner};$
\ENDIF
\IF{$v\_pin \ne p\_pin$}
    \STATE $\text{SIGNAL SQLSTATE '45000'} \implies \text{Invalid Authentication PIN};$
\ENDIF
\STATE $\text{UPDATE Tasks SET Status} = \text{'completed'} \text{ WHERE Task\_ID} = p\_task\_id;$
\STATE $\text{UPDATE Escrow\_Transactions SET Escrow\_State} = \text{'RELEASED'},$
\STATE \quad $\text{Released\_At} = \text{CURRENT\_TIMESTAMP WHERE Task\_ID} = p\_task\_id;$
\STATE $\text{UPDATE Students SET Balance} = \text{Balance} + v\_bounty$
\STATE \quad $\text{WHERE Reg\_No} = p\_runner\_id;$
\STATE $\text{COMMIT};$
\end{algorithmic}
\end{algorithm}

Because line~4 acquires an exclusive row-level lock on the target task, any concurrent payout execution blocks until the active transaction commits or rolls back. When the blocked thread acquires the lock, it immediately reads $v\_status = \text{'completed'}$, failing the precondition at line~5 and aborting cleanly without modifying balances.

\section{Event-Driven Active Database Integrity}
To enforce system invariants independently of application-level bugs, we embed business logic within database triggers.

\subsection{Self-Bidding Prevention}
Marketplace integrity requires that task requesters cannot place bids on their own errands. The trigger \texttt{trg\_prevent\_self\_bid} fires before insert operations on the \texttt{Bids} table:
\begin{align*}
    &\text{\textbf{TRIGGER} trg\_prevent\_self\_bid \textbf{BEFORE INSERT ON} Bids} \\
    &\text{\textbf{BEGIN}} \\
    &\quad \text{SELECT Requester\_ID INTO } v\_req \text{ FROM Tasks WHERE Task\_ID} = \text{NEW.Task\_ID;} \\
    &\quad \text{\textbf{IF} } v\_req = \text{NEW.Runner\_ID \textbf{THEN}} \\
    &\quad\quad \text{SIGNAL SQLSTATE '45000'} \text{ SET MESSAGE\_TEXT} = \text{'Requester cannot bid on own task';} \\
    &\quad \text{\textbf{END IF;}} \\
    &\text{\textbf{END}}
\end{align*}

\subsection{Automated Bid Resolution}
When a requester assigns a runner to an open errand, competing bids must be resolved atomically. The trigger \texttt{trg\_after\_task\_assignment} executes after updates on \texttt{Tasks}, transitioning the winning runner's bid to \texttt{accepted} and automatically marking all other pending bids for that task as \texttt{rejected}.

\section{Empirical Evaluation \& Experimental Results}
We conduct comprehensive experimental benchmarking to evaluate the performance gains of our composite indexing scheme and the resilience of our pessimistic concurrency control under extreme race conditions.

\subsection{Experimental Setup}
The evaluation environment is configured as follows:
\begin{itemize}
    \item \textbf{Database Engine:} MySQL 8.0.41 Community Server running InnoDB storage engine.
    \item \textbf{Configuration:} Transaction isolation level set to \texttt{REPEATABLE READ}, connection pool size = 10, maximum overflow = 20, \texttt{pool\_pre\_ping = True}.
    \item \textbf{Application Runtime:} FastAPI 0.146.0 running on Uvicorn with standard ASGI workers and Python 3.14.
    \item \textbf{Scale Dataset:} 10,000 synthetic tasks and bids generated via \texttt{seed\_scale\_data.py} reflecting realistic status distributions (85\% completed, 10\% cancelled, 5\% open).
\end{itemize}

\subsection{Index Performance Benchmark}
We evaluate the composite B-Tree index \texttt{idx\_tasks\_status\_id} on \texttt{Tasks(Status, Task\_ID)} and \texttt{idx\_bids\_runner\_status} on \texttt{Bids(Runner\_ID, Bid\_Status)}. Each query is sampled across 500 timed iterations preceded by 20 warm-up passes to eliminate disk I/O caching artifacts.

\begin{table*}[t]
\centering
\caption{Empirical Query Execution Plan Analysis and Runtime Latency Benchmark ($N = 500$ Iterations)}
\label{tab:index_benchmark}
\begin{tabular}{lllrrrr}
\toprule
\textbf{Target Workload} & \textbf{Access Path} & \textbf{Access Mechanism} & \textbf{Est. Rows Scanned} & \textbf{Mean Latency} & \textbf{P95 Latency} & \textbf{Speedup / Gain} \\
\midrule
Marketplace Feed & Full Table Scan & Engine Scan (\texttt{ALL}) & 9,891 & 6.8837~ms & 8.3103~ms & Baseline \\
Marketplace Feed & Indexed & \texttt{idx\_tasks\_status\_id} & 495 & 4.6063~ms & 5.4223~ms & \textbf{+33.1\% faster} \\
\midrule
Runner Bids Query & Full Table Scan & Engine Scan (\texttt{ALL}) & $\approx 10,000$ & 4.7457~ms & --- & Baseline \\
Runner Bids Query & Indexed & \texttt{idx\_bids\_runner\_status} & $\approx 5$ & 0.4226~ms & --- & \textbf{+91.1\% ($11.2\times$ faster)} \\
\bottomrule
\end{tabular}
\end{table*}

As detailed in Table~\ref{tab:index_benchmark}, the indexed access path for the marketplace feed reduces the row scan volume from 9,891 to 495 rows, delivering a 33.1\% reduction in mean latency (from 6.8837~ms to 4.6063~ms) and reducing 95th-percentile tail latency from 8.3103~ms to 5.4223~ms. For runner bid lookups, the composite key reduces execution time from 4.7457~ms to 0.4226~ms, representing an $11.2\times$ acceleration due to direct B-Tree index traversal.

\subsection{Concurrency & Race Condition Analysis}
To test the pessimistic locking engine, we deploy two multi-threaded contention benchmarks.

\subsubsection{Two-Thread Parallel Payout Race}
In the first test, two concurrent threads (\texttt{Thread-Alpha} and \texttt{Thread-Beta}) simultaneously submit the valid 4-digit PIN for the same assigned errand. As recorded in Table~\ref{tab:two_thread_race}:
\begin{table}[htbp]
\centering
\caption{Concurrent Payout Execution Under Exclusive Row Lock}
\label{tab:two_thread_race}
\begin{tabular}{lcrl}
\toprule
\textbf{Execution Thread} & \textbf{HTTP Code} & \textbf{Latency} & \textbf{Server Response} \\
\midrule
Thread-Alpha & 200 OK & 14.97~ms & PIN confirmed! Escrow released. \\
Thread-Beta  & 400 Bad Request & 15.09~ms & Invalid State: Task not active. \\
\bottomrule
\end{tabular}
\end{table}

\texttt{Thread-Alpha} acquired the row lock first and completed the transaction in 14.97~ms. \texttt{Thread-Beta} was forced to wait on the lock for 15.09~ms, upon which the task status had transitioned to \texttt{completed}, causing an immediate rejection without secondary escrow release.

\subsubsection{High-Volume 10-Thread Double-Spend Attack}
We further subjected the handover procedure to an intensive 10-worker concurrent stress test firing simultaneous payout claims.

\begin{table}[htbp]
\centering
\caption{High-Volume 10-Thread Concurrent Double-Spend Stress Test}
\label{tab:ten_thread_stress}
\begin{tabular}{lcrl}
\toprule
\textbf{Thread ID} & \textbf{HTTP Code} & \textbf{Latency} & \textbf{Server Outcome} \\
\midrule
Thread \#3 & 200 OK & 29.80~ms & Payout Successful \\
Thread \#1 & 400 Bad Request & 30.50~ms & Task not active \\
Thread \#2 & 400 Bad Request & 31.23~ms & Task not active \\
Thread \#4 & 400 Bad Request & 29.82~ms & Task not active \\
Thread \#5 & 400 Bad Request & 31.37~ms & Task not active \\
Thread \#6 & 400 Bad Request & 30.69~ms & Task not active \\
Thread \#7 & 400 Bad Request & 29.45~ms & Task not active \\
Thread \#8 & 400 Bad Request & 30.20~ms & Task not active \\
Thread \#9 & 400 Bad Request & 29.11~ms & Task not active \\
Thread \#10 & 400 Bad Request & 29.45~ms & Task not active \\
\bottomrule
\end{tabular}
\end{table}

As documented in Table~\ref{tab:ten_thread_stress}:
\begin{itemize}
    \item Total parallel requests: 10
    \item Successful payouts: \textbf{1 (Expected: 1)}
    \item Rejected duplicate claims: \textbf{9 (Expected: 9)}
    \item Pre-attack runner balance: Rs.~970.00
    \item Post-attack runner balance: Rs.~1030.00
    \item Net amount disbursed: \textbf{Rs.~60.00 (Exact single bounty)}
    \item Total execution time: 33.56~ms
\end{itemize}
The audit confirms that zero double-spending occurred, proving the mathematical invariant of the pessimistic lock boundary.

\section{Related Work}
Peer-to-peer delivery and spatial crowdsourcing have been widely studied in urban computing \cite{tong2020spatial}. Platforms such as TaskRabbit and Uber rely on distributed task allocation \cite{chen2019ridesharing}. While decentralized blockchain architectures offer trustless escrow \cite{wood2014ethereum}, they introduce prohibitive transaction fees and multi-second block confirmation latencies that are unsuitable for real-time campus micro-tasks. Centralized relational databases executing two-phase locking provide the required sub-20~ms transaction latencies while preserving absolute consistency \cite{gray1992transaction,bernstein1987concurrency}.

\section{Conclusion}
In this paper, we presented the architectural design, formal mathematical normalization, and empirical evaluation of an ACID-compliant peer-to-peer campus task dispatcher. By decomposing relations to Boyce-Codd Normal Form, we guaranteed lossless-join reconstruction and complete functional dependency preservation. By employing pessimistic row-level locking (\texttt{SELECT ... FOR UPDATE}) within MySQL InnoDB stored procedures and enforcing business rules via triggers, the system completely eliminated double-spending and phantom claims under high multi-threaded contention. Benchmark evaluations confirmed that composite B-Tree indexing accelerates open task retrieval by 33.1\% and runner bid queries by $11.2\times$, while resolving concurrent payouts in under 34~ms with zero financial leakage. Future work will investigate geo-spatial indexing and automated reputation decay models.

\begin{thebibliography}{99}

\bibitem{gray1992transaction}
J.~Gray and A.~Reuter, \emph{Transaction Processing: Concepts and Techniques}. Morgan Kaufmann, 1992.

\bibitem{chen2019ridesharing}
L.~Chen, D.~Zhang, L.~Wang, D.~Zhang, X.~Ma, and C.~Bettstetter, ``Dynamic ridesharing with flexible pickup and drop-off points,'' \emph{IEEE Transactions on Intelligent Transportation Systems}, vol.~20, no.~9, pp.~3582--3596, 2019.

\bibitem{codd1974further}
E.~F. Codd, ``Recent investigations into relational database systems,'' in \emph{IFIP Congress 1974}, pp.~1017--1021, 1974.

\bibitem{heath1971unacceptable}
I.~J. Heath, ``Unacceptable file operations in relational database systems,'' in \emph{Proceedings of the 1971 ACM SIGFIDET Workshop on Data Description, Access and Control}, pp.~19--35, 1971.

\bibitem{bernstein1987concurrency}
P.~A. Bernstein, V.~Hadzilacos, and N.~Goodman, \emph{Concurrency Control and Recovery in Database Systems}. Addison-Wesley, 1987.

\bibitem{tong2020spatial}
Y.~Tong, Z.~Zhou, Y.~Zeng, L.~Chen, and C.~Shahabi, ``Spatial crowdsourcing: a survey,'' \emph{The VLDB Journal}, vol.~29, no.~1, pp.~217--250, 2020.

\bibitem{wood2014ethereum}
G.~Wood, ``Ethereum: A secure decentralised generalised transaction ledger,'' \emph{Ethereum Project Yellow Paper}, vol.~151, pp.~1--32, 2014.

\end{thebibliography}

\end{document}